\section{Conclusion}\label{sec:conclusion}
This study asked how linguistic relatedness, domain match and documented pretraining exposure are associated with NER transfer among six South African Bantu languages. Relatedness is the most consistent predictor (RQ1), and it stays significant once corpus terms and the tokenizer-fertility gap are added; however, matching the target's corpus overrides it for isiZulu and for isiXhosa news. Targets seen in pretraining show no advantage (RQ2), although the design cannot separate exposure from cluster and fertility. Target data closes 57--89\% of transfer gaps with 500 sentences, and it does so fastest for weak cross-corpus sources (RQ3). For Setswana, Swahili is a stronger source than Sepedi or Sesotho, also with equal entity counts (RQ4).

Since the two corpora differ in both text and annotation guidelines, the corpus effects in RQ1 and RQ4 are associations whose cause remains open. Therefore, the most direct next question is whether a Sotho-Tswana source annotated consistently with the target would overtake Swahili, and the experiments in Section~\ref{sec:future} are designed to answer it.
